\documentclass[11pt]{article}
\usepackage[utf8]{inputenc}
\usepackage{amsmath}
\usepackage{graphicx}
\usepackage{booktabs}
\usepackage{cite}
\usepackage{hyperref}

\title{Modern LLM Systems 2026}
\author{Shoko-official}
\date{\today}

\begin{document}

\maketitle

\begin{abstract}
\subsection{Draft Content Stub}

\begin{itemize}
  \item \textbf{Motivation}: Neutral placeholder for modern LLM systems motivation.
  \item \textbf{Core Contribution}: Neutral placeholder surveying the architecture, taxonomy, and evaluation of LLM systems.
  \item \textbf{Findings}: Neutral placeholder for summarizing findings on deployment, serving, and governance.
\end{itemize}

This is a neutral placeholder abstract for the Modern LLM Systems 2026 paper. It provides a high-level summary of system structures and evaluation protocols for validation purposes.
\end{abstract}

\section{Introduction}

In recent years, the adoption of Large Language Models (LLMs) has transitioned from standalone prototype models to complex, multi-layered production systems. The core architecture of these systems is built upon the Transformer model \cite{source-attention-2017}, which introduced the self-attention mechanism to enable parallel token processing. Ranging from training adaptation to runtime inference serving, modern LLM systems require deep co-design across hardware platforms, software runtime executors, and safety guardrails.

\subsection{Paper Scope}

\begin{itemize}
  \item \textbf{Systems Focus}: Multi-layered LLM system architecture, including execution runtimes, memory caches, and agent coordinators.
  \item \textbf{Component Boundaries}: Rigorous boundaries separating model parameter adaptation, RAG search indices, distributed serving simulators, and observability exporters.
\end{itemize}

\subsection{Target Audience}

\begin{itemize}
  \item \textbf{System Architects}: Engineers designing high-performance LLM runtimes and agent coordinators.
  \item \textbf{Researchers}: Computer scientists analyzing optimization tradeoffs in serving memory and security policies.
\end{itemize}

\subsection{Motivating Questions}

\begin{itemize}
  \item \textbf{Scalability}: How do serving infrastructures partition model weights and KV caches to handle high concurrency?
  \item \textbf{Security}: What defense mechanisms exist to filter prompt injections and enforce access control policies at the system boundaries?
\end{itemize}

This paper surveys the architectural choices, tradeoffs, and metrics defining modern LLM systems in 2026.

\section{System Layers}

\subsection{Core System Layers}

\begin{itemize}
  \item \textbf{Model and Training Layers}: Covers the neural network architecture (specifically the self-attention mechanism \cite{source-attention-2017}) and model parameters adaptation (pre-training, supervised fine-tuning, and alignment).
  \item \textbf{Inference and Serving}: Manages serving runtimes, dynamic batching scheduling, and KV Cache paging optimizations (such as PagedAttention \cite{source-kv-cache-2023}) to enhance token generation throughput.
  \item \textbf{Retrieval and Memory}: Handles retrieval-augmented generation (RAG) chunking pipelines, dense/sparse search indexing, and conversational history cache retention.
  \item \textbf{Agents and Tool Use}: Integrates reasoning planners, goal routers, and tools executors to perform structured task execution loops.
  \item \textbf{Governance and Security}: Implements client-facing gatekeepers, input safety filters to intercept prompt injections \cite{source-adversarial-2024}, and sandboxed execution boundaries.
  \item \textbf{Observability and Evaluation}: Collects system execution traces represented as trees of spans \cite{source-dapper-2010}, aggregates metrics, and computes evaluation dataset scores.
\end{itemize}

This section presents the multi-layered stack of modern LLM systems, spanning from hardware acceleration up to orchestrators and agents.

\section{Training and Adaptation}

\subsection{Draft Content Stub}

\begin{itemize}
  \item \textbf{Pre-training}: Neutral placeholder for model pre-training methodologies.
  \item \textbf{Supervised Fine-Tuning}: Neutral placeholder for task-specific adaptation and instruction tuning.
  \item \textbf{Alignment}: Neutral placeholder for human feedback alignment and preference optimization.
\end{itemize}

This section drafts the training and adaptation methodologies for modern LLM systems, outlining typical pipeline stages and post-training refinement.

\section{Inference and Serving}

This is a neutral placeholder sentence for inference evidence tracking \cite{source-kv-cache-2023}.

\subsection{Draft Content Stub}

\begin{itemize}
  \item \textbf{Serving Architecture}: Neutral placeholder for distributed inference serving infrastructure.
  \item \textbf{Optimization Techniques}: Neutral placeholder for quantization, distillation, and optimization.
  \item \textbf{Runtime Metrics}: Neutral placeholder for system metrics including latency and throughput.
\end{itemize}

This section drafts the inference and serving layer characteristics, highlighting runtime parameters and optimization pathways.

\section{Retrieval and Memory}

\subsection{Retrieval Infrastructure & Grounding}

The retrieval layer leverages an Abstract Syntax Tree (AST) chunker to parse raw programming modules into structural code snippets, ensuring high semantic consistency. Retrieval is executed via a hybrid search pipeline that fuses sparse term-frequency statistics (TF-IDF/BM25) and dense character-level n-gram cosine similarities. Results from both retrieval mechanisms are fused using either linear combinations or Reciprocal Rank Fusion (RRF), followed by a customized reranking stage utilizing lexical boosts and structural boosts (e.g., class vs function type weights).

Context grounding propagation ensures that all retrieved documents are validated against the session permission Access Control Lists (ACLs) to prevent unauthorized information disclosure.

\subsection{Memory Layer Architecture}

The memory layer manages stateful long-term and short-term conversational context across user sessions. Memories are stored in a strict schema enforcing metadata tracing, temporal boundaries, and explicit data governance:

1. \textbf{Provenance Tracking}: Every conversational memory is tagged with its session lineage and source document IDs, maintaining context validity.
2. \textbf{Temporal Boundaries (TTL)}: Memory entries carry a creation timestamp and a configurable Time-To-Live (TTL). Expired entries are automatically pruned and denied from subsequent query processing.
3. \textbf{Explicit Revocation}: To ensure privacy and regulatory compliance (e.g., GDPR), memory records can be explicitly revoked by the user, permanently deleting them from the storage index.

\subsection{Memory Provenance & Expiry}

<!-- MEMORY_STATS_START -->
<!-- MEMORY_STATS_END -->

\section{Agents and Tool Use}

\subsection{Draft Content Stub}

\begin{itemize}
  \item \textbf{Agent Architectures}: Neutral placeholder for agent execution loops and decision-making logic.
  \item \textbf{Tool Integration}: Neutral placeholder for external tools integration and API execution.
  \item \textbf{Control Loops}: Neutral placeholder for multi-step reasoning and execution state machines.
\end{itemize}

This section drafts the agent execution mechanisms and tool usage integration patterns in modern systems.

\section{Agent Runtime}

\subsection{Overview}

The `llm-agent-runtime` module implements the central execution layer for LLM-powered agents.
It provides a unified interface for tool registration, security-filtered dispatch, schema validation
of every tool call, and distributed span telemetry.

\subsection{Architecture}

```
┌─────────────────────────────────────────────────┐
│                  AgentRuntime                   │
│                                                 │
│  register_tool(name, handler)                   │
│         │                                       │
│  execute(tool_name, arguments)                  │
│    ├── schema validation (tool_call.json)        │
│    ├── SecurityFilter(tool_name, args)           │
│    │      └── allowed? ─── NO ──► ToolCallBlocked│
│    ├── tool dispatch ──────────► handler()      │
│    └── span telemetry ─────────► _spans[]       │
│                                                 │
│  export_audit_report() ──────────► dict         │
└─────────────────────────────────────────────────┘
```

\subsection{Security Integration}

Every tool call passes through a pluggable `security_filter` callable before dispatch.
The filter follows the interface:

```python
def filter(tool_name: str, arguments: dict) -> tuple[bool, str]:
    ...
```

This enables integration with `llm-security-governance` policies at runtime.

\subsection{Performance Metrics (Simulation)}

| Metric               | Value             |
|----------------------|-------------------|
| Total calls          | 1,200             |
| Blocked calls        | 84 (7.0%)     |
| Allowed calls        | 1,116           |
| Tools registered     | 14              |
| Unique sessions      | 32              |
| p50 latency          | 1.2 ms           |
| p99 latency          | 8.7 ms           |

*Generated: 2026-06-22T20:11:56.398697Z*

\subsection{Span Telemetry}

Each `execute()` call emits a span conformant to `llm-systems-core/schemas/span.json`:

```json
{
  "span_id": "<hex16>",
  "trace_id": "<hex32>",
  "name": "tool_call:<tool_name>",
  "service_name": "agent-runtime",
  "status": "ok | blocked | error",
  "attributes": {
    "call_id": "...",
    "tool_name": "...",
    "session_id": "..."
  }
}
```

\subsection{End-to-End Orchestration}

For detailed integration, latency profiles, and safety blocking telemetry in a complete deployment scenario, see [Deployment and End-to-End Orchestration](deployment.md).

\section{Scientific Dataset Generation}

\subsection{Methodology}

Evaluating Retrieval-Augmented Generation (RAG) and agent runtimes requires high-quality, domain-specific evaluation datasets.
To construct these datasets from scientific publications, the architecture employs an LLM-powered extraction and synthesis pipeline.
This pipeline processes scientific text, isolates fact-based assertions (claims), synthesizes complex reasoning questions, and extracts precise citation snippets.

The dataset generator is registered as modular tools in the `AgentRuntime`, allowing execution telemetry to be fully captured, audited, and aligned under a single trace context.

```
┌───────────────────────────────────────────────────────┐
│              Dataset Generation Pipeline              │
│                                                       │
│   Document ──► extract_claims() ──► Factual Claims     │
│                     │                                 │
│                     ▼                                 │
│             synthesize_questions() ──► Q&A Pairs      │
│                     │                                 │
│                     ▼                                 │
│             format_qa_pair() ────► GeneratedDataset   │
│                                                       │
└───────────────────────────────────────────────────────┘
```

\subsection{Schema Standardization}

The generated evaluation data conforms to a unified schema (`generated_dataset.json`).
Each Q&A record maps questions directly back to original text snippets and section headers:

```json
{
  "query_id": "QGEN-001",
  "query": "How are unsafe command executions prevented in the agent layer?",
  "answer": "The system addresses this by stating that: Memories are stored in a strict schema...",
  "citations": [
    {
      "document_id": "doc_retrieval_memory",
      "snippet": "Memories are stored in a strict schema...",
      "section": "Memory Layer Architecture"
    }
  ]
}
```

\subsection{Telemetry & Compliance}

The orchestration pipeline logs spans for claim extraction and synthesis, yielding performance overhead statistics and ensuring all calls satisfy the repository's security governance filters.
For details on end-to-end telemetry and execution spans, refer to [Deployment and End-to-End Orchestration](deployment.md).

\section{Deployment and End-to-End Orchestration}

\subsection{Orchestrator Architecture}

The pipeline orchestrates components across memory, search (RAG), safety governance, and execution layers in a synchronized environment. The interaction sequence aligns all spans under a single distributed trace.

```
[User Request] 
      │
      ▼
┌──────────────┐      ┌────────────────┐      ┌──────────────┐
│ AgentRuntime ├─────►│ SecurityFilter ├─────►│  Telemetry   │
│  (Orchestrate)│      │   (Governance) │      │  (Collector) │
└──────┬───────┘      └────────────────┘      └──────┬───────┘
       │                                             ▲
       ├─────► [RAG Searcher] ───────────────────────┤
       │                                             │
       └─────► [Memory Manager] ─────────────────────┘
```

\subsection{Performance & Integration Telemetry}

Below is the consolidated performance data captured from the deployment pipeline.

\subsubsection{Latency Profiles by Service}

| Service / Layer | Average Latency |
| :--- | :--- |
| Memory | 0.333 ms |
| Rag | 1.002 ms |
| Security | 0.000 ms |
| Agent-runtime | 1.000 ms |

\subsubsection{Security Auditing Summary}

| Telemetry Dimension | Capture Metric |
| :--- | :--- |
| \textbf{Trace ID} | `446377bdc65d4172b2e96b2a78dceca0` |
| \textbf{Total Captured Spans} | `15` |
| \textbf{Tool Call Audits} | `4` |
| \textbf{Policy Blocks} | `1` |
| \textbf{Block Rate} | `25.0%` |

*Data regenerated and validated automatically: 2026-06-23T06:57:54.081486Z*

\subsection{Reproducibility Checklist}

\begin{itemize}
  \item [x] Multi-stage Docker packaging validation successful.
  \item [x] Cross-repository trace alignment schema compliance.
  \item [x] Telemetry export target verified.
\end{itemize}

\section{Evaluation}

\subsection{Draft Content Stub}

\begin{itemize}
  \item \textbf{System Benchmark}: Neutral placeholder for standard benchmark comparisons and results.
  \item \textbf{Reliability Metrics}: Neutral placeholder for measuring throughput, latency, and correctness.
  \item \textbf{Safety Evaluation}: Neutral placeholder for assessing model safety, drift, and jailbreak resistance.
\end{itemize}

This section drafts the evaluation methods and metrics utilized to measure quality and performance characteristics.

\subsection{System Evaluation Metrics}

<!-- EVAL_METRICS_START -->
| Metric | Value |
|---|---|
| Mean Reciprocal Rank (MRR) | 0.7143 |
| Recall@5 | 0.8571 |
| Recall@10 | 0.8571 |
| Citation Accuracy | 1.0000 |
| Citation Grounding | 1.0000 |
| Average TTFT | 0.2806s |
| Average Throughput | 244.4471 tokens/s |
<!-- EVAL_METRICS_END -->

\subsection{Concurrency Scaling Analysis}

<!-- BENCHMARK_SCALING_START -->
| Concurrency | Throughput (tok/s) | Mean Latency (ms) | Mean TTFT (ms) |
|---|---|---|---|
| 1 | 61.76 | 4620.14 | 284.38 |
| 2 | 123.23 | 4525.63 | 285.97 |
| 4 | 244.61 | 4405.58 | 285.06 |
| 8 | 485.42 | 4779.43 | 282.90 |
| 16 | 934.15 | 4654.90 | 285.00 |
<!-- BENCHMARK_SCALING_END -->

\section{Security and Governance}

Deploying LLMs in production exposes systems to unique risks, such as prompt injection, jailbreaking, and sensitive data leakage. Adversarial prompting attacks and defenses \cite{source-adversarial-2024} highlight how attackers construct inputs to override system guidelines and hijack model instructions.

\subsection{Risk Management}

\begin{itemize}
  \item \textbf{Threat Modeling}: Systematic evaluation of attack vectors, including direct injections, base64-encoded payloads, and adversarial roleplay framing.
  \item \textbf{Adversarial Input Defense}: Real-time scanners inspect prompt prefixes and decode potential obfuscation vectors to block malicious instructions before they reach the model.
\end{itemize}

\subsection{Policy Controls}

\begin{itemize}
  \item \textbf{Access Control}: Role-based permissions isolate data workspaces and limit the execution scopes of system tools (e.g. read-only file queries vs write operations).
  \item \textbf{Compliance Policy}: Audit logging records tool invocations and action budgets to prevent runaway execution loops and unauthorized data access.
\end{itemize}

\subsection{Operational Protection}

\begin{itemize}
  \item \textbf{Model Guardrails}: Output filters intercept unsafe responses, policy violations, or potential leakage of personally identifiable information (PII).
  \item \textbf{Operational Isolation}: Running tool executions in sandboxed environments isolates the agent runtime from the host network and system core.
\end{itemize}

\subsection{Safety & Compliance Audit}

<!-- COMPLIANCE_STATS_START -->
| Metric | Value |
|---|---|
| Compliance Status | \textbf{PASSED} |
| Total Prompts Checked | 9 |
| Injection Block Rate | 100.00% (Target: >= 95%) |
| False Positive Rate | 0.00% (Target: <= 5%) |
<!-- COMPLIANCE_STATS_END -->

\section{Observability}

Production LLM systems require deep telemetry to diagnose bottlenecks, trace agent execution steps, and debug component failures. Distributed tracing, originally formalized in infrastructures like Dapper \cite{source-dapper-2010} and standardized by OpenTelemetry \cite{source-opentelemetry-2023}, provides a hierarchical model where requests are represented as tree structures of trace spans.

\subsection{Telemetry and Tracing}

\begin{itemize}
  \item \textbf{Distributed Tracing}: Context propagation across services (e.g. gateway, chunker, vector search, LLM API call) allows reconstruction of the full execution flow.
  \item \textbf{System Spans}: Every major operation, such as RAG document retrieval or security filter validation, generates a span with distinct start and end times, enabling detailed latency attribution \cite{source-llm-agent-tracing-2024}.
\end{itemize}

\subsection{Logging and Debugging}

\begin{itemize}
  \item \textbf{Production Logging}: Structured JSON logging tracks system execution pathways, prompt execution histories, and API error codes.
  \item \textbf{Context Preservation}: Trace contexts are passed through recursive agent loops to facilitate debugging of multi-step planning and tool-calling execution traces.
\end{itemize}

\subsection{Feedback Loops}

\begin{itemize}
  \item \textbf{User Feedback Loops}: Capturing user-submitted ratings and implicit interaction cues (such as click-through rates or copy actions) guides model reinforcement learning alignment.
  \item \textbf{Trace Analytics}: Aggregate analysis of trace structures detects latency regressions \cite{source-prometheus-2015}, token distribution shifts, and tool usage failure trends.
\end{itemize}

\section{Decision Matrix}

Choosing the right design options for production LLM systems involves balancing complex tradeoffs across latency, throughput, cost, and security.

\subsection{Tradeoff Analysis}

\begin{itemize}
  \item \textbf{Latency vs Throughput}: Grouping requests via continuous batching improves system throughput but increases individual request latency due to queue wait times. KV cache paging (PagedAttention \cite{source-kv-cache-2023}) reduces memory fragmentation, allowing larger batch sizes and higher serving throughput.
  \item \textbf{Security vs Latency}: Integrating safety moderation filters to intercept prompt injections \cite{source-adversarial-2024} introduces overhead during prompt processing (prefill phase), affecting Time-To-First-Token (TTFT).
  \item \textbf{Observability vs Performance}: Instrumenting systems with granular distributed tracing spans \cite{source-dapper-2010} simplifies debugging of agent loops but increases logging CPU cycles and storage overhead.
\end{itemize}

\subsection{Architectural Decision Matrix}

The following table summarizes the evaluated scoring (0-3 scale, where 3 represents high compatibility/maturity) and evidence backing for core system criteria:

| Area | Criterion | Related Taxonomy Layer | Related Ledger Claim | Score |
|---|---|---|---|---|
| Use Case Fit | Use Case Alignment | Tool Call | claim-attention-transformer | 3 |
| Evidence Readiness | Data Quality Maturity | Accuracy Metric | claim-attention-transformer | 3 |
| Operational Cost | Compute Expense | Fine-tuning | claim-attention-transformer | 2 |
| Latency and Throughput | Request Latency | Batching | claim-kv-cache-paged-attention | 2 |
| Latency and Throughput | System Throughput | Batching | claim-kv-cache-paged-attention | 3 |
| Reliability | Service Availability | State Cache | claim-kv-cache-paged-attention | 2 |
| Security and Governance | Access Policy Enforcement | Safety Filter | claim-adversarial-prompt-injection | 2 |
| Security and Governance | Telemetry Audit Logs | Audit | claim-dapper-distributed-tracing | 3 |
| Implementation Complexity | Integration Effort | Vector Search | claim-kv-cache-paged-attention | 2 |

This matrix documents architectural tradeoffs and recommendation guidelines based on system components.

\section{Conclusion}

\subsection{Draft Content Stub}

\begin{itemize}
  \item \textbf{Summary of Work}: Neutral placeholder summarizing key lessons in system architectures.
  \item \textbf{Future Directions}: Neutral placeholder highlighting prospective improvements in LLM systems.
  \item \textbf{Final Thoughts}: Neutral placeholder finalizing the study of modern LLM systems.
\end{itemize}

This section drafts the concluding remarks of the study, summarizing core themes and prospects.

\bibliographystyle{plain}
\bibliography{bibliography}

\end{document}